\begin{helper}[Auxiliary colors in the no-repetition case]
Let \((\varphi,T)\) be a minimum bad pair with
\(2<q:=\rho(T)<n:=|E(T)|\), satisfying the no-repetition condition on
\(y_q,\ldots,y_n\), and let \(\beta\) be missing at \(y_n\).  There
exist colors \(\gamma,\varepsilon\) and an index \(m\) such that
\(\gamma\ne\beta\), \(\varepsilon\ne\gamma\), \(\gamma\) is missing
at \(y_q\), \(\varepsilon\) is missing at \(y_m\),
\(m+2\le q\), and both \(\gamma\) and \(\varepsilon\) are unused
through \(y_q\).
\end{helper}
\begin{proof}
The \(k\)-triple clause of \noderef{TashTreeBadPair} gives
\(k\ge\Delta(G)+1\).  Therefore
\noderef{MissingColorsNonemptyOfMaximumDegree} supplies a color
\(\gamma\) missing at \(y_q\).

This color is unused through \(y_q\).  If a formal edge index \(r\)
satisfied \(1\le r<q\) and the corresponding exact edge had color
\(\gamma\), then \noderef{TashkinovTree} would supply a witness
\(y_w\) with \(w\le r<q\) at which \(\gamma\) is missing.  Both
\(y_w\) and the distinct vertex \(y_q\) belong to the proper prefix
\(T_q\), contradicting its elementariness from
\noderef{TashTreeProperPrefixElementary}.  This is exactly
\noderef{TreeColorUnusedThrough}; in particular, as required by its
formal definition, no condition has been imposed on the first displayed
edge at formal index zero.  Also \(\gamma\ne\beta\): otherwise the
color \(\beta\), already missing at \(y_n\), would be missing at the
suffix vertex \(y_q\), contrary to the no-repetition hypothesis.

Apply \noderef{TashTreeFourUnusedColors} to the positive proper prefix
\(T_{q-2}\).  Its at least four members of
\noderef{TreeFreeColors} are missing at vertices of that prefix by
\noderef{TreeMissingColors} and absent from its exact colored edges by
\noderef{TreeUsedColors}.  Extending through \(y_q\) introduces only
the two exact edges \(e_{q-1}\) and \(e_q\).  Each excludes at most its
single assigned color from being unused, so at least two of the four
colors remain unused through \(y_q\).  Choose among them a color
\(\varepsilon\ne\gamma\).  Its membership in
\noderef{TreeMissingColors} yields a witness \(y_m\) in
\(T_{q-2}\), hence \(m+2\le q\); its avoidance of the old prefix and
of the two additional exact edges is precisely
\noderef{TreeColorUnusedThrough} through \(y_q\).  These choices give
all the claimed properties.
\end{proof}
